% GENERATED FILE. Edit canonical lessons, then run npm run generate:books.
% canonical-source-hash: fnv1a64:429dbf54d5eb88c7
% canonical-lessons: ZH-C424-suidao, ZH-C424-diaotou, ZH-C424-lubiao, ZH-C424-shisu, ZH-C424-banlu, ZH-R424-review-420-424

\chapter{On the Road}
\label{ch:zh-on-the-road}
\hlchaptermodality{424}

\begin{chapteropening}
\textbf{By the end of this chapter:} \emph{I can talk about tunnels, turning round, road signs, speed per hour and stopping halfway.}

Complete the last lesson of chapter 424: I can talk about tunnels, turning round, road signs, speed per hour and stopping halfway.
\end{chapteropening}

\section[suìdào]{suìdào — tunnel}

\label{lesson:ZH-C424-suidao}



\begin{quote}
Before the new one: say the Mandarin for to head down south, then the Mandarin for north and south.
\end{quote}

\subsection*{You'll want to know: suìdào}
\emph{suìdào} — \textquotedblleft{}tunnel\textquotedblright{}.

\begin{quote}
\emph{qiánbian yǒu yì tiáo suìdào} — There's a tunnel up ahead.
\end{quote}

Its characters come later; for now it is a word for the ear.

\begin{cousinweb}
\emph{suìdào} is built from \emph{suì} 'a passage' and \emph{dào} 'a way', both inherited from Middle Chinese, \emph{zwijH} and \emph{dawX} (Guangyun, in Baxter's transcription), which glosses the first as the passage into a tomb and the second 'a road'.
\end{cousinweb}

\begin{grammarlens}[title={What you've built}]
One more word for getting around.
\end{grammarlens}

\subsection*{Your turn}
\begin{itemize}
\raggedright
  \item \emph{Say it:} \emph{suìdào}
  \item \emph{Say it:} \emph{suìdào}, once more
  \item \emph{Say it:} \emph{qiánbian yǒu yì tiáo suìdào} — There's a tunnel up ahead.
  \item \emph{Recall:} say the Mandarin for to head down south, then the Mandarin for north and south, then say \emph{suìdào} again
\end{itemize}

\begin{tcolorbox}[breakable,colback=teal!4,colframe=teal!35!black,title={Before you move on}]
What does suìdào mean? (\textquotedblleft{}Tunnel\textquotedblright{}.) Say it once more.
\end{tcolorbox}
\section[diàotóu]{diàotóu — to turn round (in a car); to make a U-turn}

\label{lesson:ZH-C424-diaotou}



\begin{quote}
Before the new one: say the Mandarin for north and south, then the Mandarin for a tunnel through a hill.
\end{quote}

\subsection*{You'll want to know: diàotóu}
\emph{diàotóu} — \textquotedblleft{}to turn round (in a car); to make a U-turn\textquotedblright{}.

\begin{quote}
\emph{zhèr bù néng diàotóu wǎng qián kāi} — You can't turn round here: drive on.
\end{quote}

Its characters come later; for now it is a word for the ear.

\begin{cousinweb}
\emph{diàotóu} is built from \emph{diào} 'to turn' and \zh{头} 'the head', both inherited from Middle Chinese, \emph{dewH} and \emph{duw} (Guangyun, in Baxter's transcription), which glosses the first 'to shake'; the sense 'to turn' is not given there.
\end{cousinweb}

\begin{grammarlens}[title={What you've built}]
One more word for getting around.
\end{grammarlens}

\subsection*{Your turn}
\begin{itemize}
\raggedright
  \item \emph{Say it:} \emph{diàotóu}
  \item \emph{Say it:} \emph{diàotóu}, once more
  \item \emph{Say it:} \emph{zhèr bù néng diàotóu wǎng qián kāi} — You can't turn round here: drive on.
  \item \emph{Recall:} say the Mandarin for north and south, then the Mandarin for a tunnel through a hill, then say \emph{diàotóu} again
\end{itemize}

\begin{tcolorbox}[breakable,colback=teal!4,colframe=teal!35!black,title={Before you move on}]
What does diàotóu mean? (\textquotedblleft{}To turn round (in a car); to make a U-turn\textquotedblright{}.) Say it once more.
\end{tcolorbox}
\section[lùbiāo]{lùbiāo — road sign}

\label{lesson:ZH-C424-lubiao}



\begin{quote}
Before the new one: say the Mandarin for a tunnel through a hill, then the Mandarin for to make a U-turn.
\end{quote}

\subsection*{You'll want to know: lùbiāo}
\emph{lùbiāo} — \textquotedblleft{}road sign\textquotedblright{}.

\begin{quote}
\emph{lùbiāo shuō wǎng qián zǒu} — The sign says go straight on.
\end{quote}

Its characters come later; for now it is a word for the ear.

\begin{cousinweb}
\emph{lùbiāo} is built from \zh{路} 'a road' and \emph{biāo} 'a mark', both inherited from Middle Chinese, \emph{luH} and \emph{pjiew} (Guangyun, in Baxter's transcription), which glosses the second as the tip of a tree and 'to raise'.
\end{cousinweb}

\begin{grammarlens}[title={What you've built}]
One more word for getting around.
\end{grammarlens}

\subsection*{Your turn}
\begin{itemize}
\raggedright
  \item \emph{Say it:} \emph{lùbiāo}
  \item \emph{Say it:} \emph{lùbiāo}, once more
  \item \emph{Say it:} \emph{lùbiāo shuō wǎng qián zǒu} — The sign says go straight on.
  \item \emph{Recall:} say the Mandarin for a tunnel through a hill, then the Mandarin for to make a U-turn, then say \emph{lùbiāo} again
\end{itemize}

\begin{tcolorbox}[breakable,colback=teal!4,colframe=teal!35!black,title={Before you move on}]
What does lùbiāo mean? (\textquotedblleft{}Road sign\textquotedblright{}.) Say it once more.
\end{tcolorbox}
\section[\texorpdfstring{\zh{时速} (shísù)}{时速 (shísù)}]{\zh{时速} (shísù) — speed per hour}

\label{lesson:ZH-C424-shisu}



\begin{quote}
Before the new one: say the Mandarin for to make a U-turn, then the Mandarin for a road sign.
\end{quote}

\subsection*{You'll want to know: \zh{时速}}
\textbf{\zh{时速}} — \emph{shísù} — \textquotedblleft{}speed per hour\textquotedblright{}.

\begin{quote}
\textbf{\zh{我们的时速是一百公里}} \emph{wǒmen de shísù shì yìbǎi gōnglǐ} — We're doing a hundred kilometres an hour.
\end{quote}

\begin{cousinweb}
\textbf{\zh{时速}} (\emph{shísù}) is built from \zh{时} 'an hour' and \zh{速} 'speed', both inherited from Middle Chinese, \emph{dzyi} and \emph{suwk} (Guangyun, in Baxter's transcription).
\end{cousinweb}

\begin{grammarlens}[title={What you've built}]
One more word for getting around.
\end{grammarlens}

\subsection*{Your turn}
\begin{itemize}
\raggedright
  \item \emph{Say it:} \emph{shísù}
  \item \emph{Say it:} \emph{shísù}, once more
  \item \emph{Say it:} \emph{wǒmen de shísù shì yìbǎi gōnglǐ} — We're doing a hundred kilometres an hour.
  \item \emph{Recall:} say the Mandarin for to make a U-turn, then the Mandarin for a road sign, then say \emph{shísù} again
\end{itemize}

\begin{tcolorbox}[breakable,colback=teal!4,colframe=teal!35!black,title={Before you move on}]
What does \zh{时速} mean? (\textquotedblleft{}Speed per hour\textquotedblright{}.) Say it once more.
\end{tcolorbox}
\section[\texorpdfstring{\zh{半路} (bànlù)}{半路 (bànlù)}]{\zh{半路} (bànlù) — halfway; on the way}

\label{lesson:ZH-C424-banlu}



\begin{quote}
Before the new one: say the Mandarin for a road sign, then the Mandarin for speed per hour.
\end{quote}

\subsection*{You'll want to know: \zh{半路}}
\textbf{\zh{半路}} — \emph{bànlù} — \textquotedblleft{}halfway; on the way\textquotedblright{}.

\begin{quote}
\textbf{\zh{我们在半路停了车}} \emph{wǒmen zài bànlù tíng le chē} — We stopped the car on the way.
\end{quote}

\begin{cousinweb}
\textbf{\zh{半路}} (\emph{bànlù}) is built from \zh{半} 'half' and \zh{路} 'a road', both inherited from Middle Chinese, \emph{panH} and \emph{luH} (Guangyun, in Baxter's transcription).
\end{cousinweb}

\begin{grammarlens}[title={What you've built}]
One more word for getting around.
\end{grammarlens}

\subsection*{Your turn}
\begin{itemize}
\raggedright
  \item \emph{Say it:} \emph{bànlù}
  \item \emph{Say it:} \emph{bànlù}, once more
  \item \emph{Say it:} \emph{wǒmen zài bànlù tíng le chē} — We stopped the car on the way.
  \item \emph{Recall:} say the Mandarin for a road sign, then the Mandarin for speed per hour, then say \emph{bànlù} again
\end{itemize}

\begin{tcolorbox}[breakable,colback=teal!4,colframe=teal!35!black,title={Before you move on}]
What does \zh{半路} mean? (\textquotedblleft{}Halfway; on the way\textquotedblright{}.) Say it once more.
\end{tcolorbox}
\section[(dialogue)]{Looking back: from Parking to On the Road}

\label{lesson:ZH-R424-review-420-424}



\begin{quote}
Say the Mandarin for a driving licence and a traffic lane.
\end{quote}

\subsection*{Your turn}
\begin{itemize}
\raggedright
  \item \emph{Say it:} a driving licence, a traffic lane, a parking space, to reverse the car, to break the traffic rules
  \item \emph{Say it:} a main road between towns, the expressway, the speed of a vehicle, a toll gate, to break the speed limit
  \item \emph{Say it:} to navigate by satnav, the northeast, the northwest, to head up north, high-speed
  \item \emph{Say it:} road conditions, the southeast, the southwest, to head down south, north and south
  \item \emph{Say it:} a tunnel through a hill, to make a U-turn, a road sign, speed per hour, halfway
  \item \emph{Read it:} \zh{速} \zh{北} \zh{南} — each one aloud, then a word that holds it
\end{itemize}

\begin{tcolorbox}[breakable,colback=teal!4,colframe=teal!35!black,title={Before you move on}]
A driving licence? (\emph{jiàzhào}.) A traffic lane? (\textbf{\zh{车道}}, \emph{chēdào}.) A parking space? (\emph{chēwèi}.) To reverse the car? (\emph{dàochē}.) To break the traffic rules? (\emph{wéizhāng}.) A main road between towns? (\textbf{\zh{公路}}, \emph{gōnglù}.) The expressway? (\textbf{\zh{高速}}, \emph{gāosù}.) The speed of a vehicle? (\textbf{\zh{车速}}, \emph{chēsù}.) A toll gate? (\emph{shōufèizhàn}.) To break the speed limit? (\emph{chāosù}.) To navigate by satnav? (\emph{dǎoháng}.) The northeast? (\textbf{\zh{东北}}, \emph{dōngběi}.) The northwest? (\textbf{\zh{西北}}, \emph{xīběi}.) To head up north? (\textbf{\zh{北上}}, \emph{běishàng}.) High-speed? (\textbf{\zh{快速}}, \emph{kuàisù}.) Road conditions? (\emph{lùkuàng}.) The southeast? (\textbf{\zh{东南}}, \emph{dōngnán}.) The southwest? (\textbf{\zh{西南}}, \emph{xīnán}.) To head down south? (\textbf{\zh{南下}}, \emph{nánxià}.) North and south? (\textbf{\zh{南北}}, \emph{nánběi}.) A tunnel through a hill? (\emph{suìdào}.) To make a U-turn? (\emph{diàotóu}.) A road sign? (\emph{lùbiāo}.) Speed per hour? (\textbf{\zh{时速}}, \emph{shísù}.) Halfway? (\textbf{\zh{半路}}, \emph{bànlù}.)
\end{tcolorbox}
